\section{Podsumowanie}

\begin{frame}{Wnioski}
\textbf{Główne osiągnięcia pracy:}

\begin{enumerate}
    \item \textbf{Teoretyczne}
    \begin{itemize}
        \item Analiza struktury optymalnych polityk w MDP z dyskontowaniem QH
        \item Redukcja polityk do postaci $(\mu^\star, \phi_s^\star)$
        \item Udowodnienie zbieżności algorytmów uczenia ze wzmocnieniem
    \end{itemize}
    
    \vspace{0.5em}
    \item \textbf{Praktyczne}
    \begin{itemize}
        \item Implementacja algorytmów QH Policy Evaluation i QH Q-Learning
        \item Weryfikacja na modelach testowych
        \item Demonstracja działania w problemie zarządzania zapasami
    \end{itemize}
    
    \vspace{0.5em}
    \item \textbf{Poznawcze}
    \begin{itemize}
        \item Zrozumienie mechanizmów niespójności czasowej w podejmowaniu decyzji
        \item Modelowanie zachowań z uprzedzeniem teraźniejszości
        \item Porównanie decydentów zobowiązanych, naiwnych i wyrafinowanych
    \end{itemize}
\end{enumerate}
\end{frame}

\begin{frame}{Wnioski -- Niespójność czasowa}
\textbf{Kluczowe odkrycie:}

\begin{block}{Niespójność czasowa w optymalnych politykach}
Dla $\alpha < 1$ optymalna reguła pierwszego kroku $\mu^\star$ różni się od optymalnej polityki stacjonarnej $\phi_s^\star$.
\end{block}

\vspace{0.5em}
\textbf{Implikacje:}
\begin{itemize}
    \item Decydent zobowiązany (precommitted) osiąga lepsze wyniki niż decydent naivny
    \item Planowanie z wyprzedzeniem i trzymanie się planu jest korzystne
    \item Istotne dla:
    \begin{itemize}
        \item Ekonomii behawioralnej (planowanie oszczędności, inwestycji)
        \item Psychologii uzależnień (programy terapeutyczne)
        \item Projektowania systemów wspomagania decyzji
    \end{itemize}
\end{itemize}

\vspace{0.5em}
\textbf{Przykład:} Student uczący się do egzaminu -- decydent zobowiązany tworzy plan nauki i się go trzyma, osiągając lepszy wynik niż student codziennie odkładający naukę "na jutro".
\end{frame}

\begin{frame}{Dalsze prace}
\textbf{Możliwe kierunki rozwoju:}

\begin{enumerate}
    \item \textbf{Rozszerzenie teorii}
    \begin{itemize}
        \item Analiza decydentów wyrafinowanych (sophisticated agents)
        \item Gry wieloagentowe z QH dyskontowaniem
        \item Nieskończone przestrzenie stanów i akcji
    \end{itemize}
    
    \vspace{0.5em}
    \item \textbf{Zaawansowane algorytmy}
    \begin{itemize}
        \item Deep Q-Learning dla dużych przestrzeni stanów
        \item Actor-Critic dla ciągłych przestrzeni akcji
        \item Transfer learning między różnymi wartościami $\alpha$
    \end{itemize}
    
    \vspace{0.5em}
    \item \textbf{Zastosowania praktyczne}
    \begin{itemize}
        \item Systemy rekomendacji uwzględniające preferencje czasowe użytkowników
        \item Planowanie finansowe osobiste (oszczędności, kredyty)
        \item Optymalizacja łańcucha dostaw z uwzględnieniem niespójności czasowej
        \item Projektowanie mechanizmów motywacyjnych (commitment devices)
    \end{itemize}
\end{enumerate}
\end{frame}

\begin{frame}{Dalsze prace (cd.)}
\textbf{Wyzwania badawcze:}

\begin{itemize}
    \item \textbf{Estymacja parametrów:} Jak wyznaczać $\alpha$ i $\beta$ z danych rzeczywistych?
    \item \textbf{Adaptacja online:} Czy parametry QH mogą się zmieniać w czasie?
    \item \textbf{Uczenie w środowisku częściowo obserwowalnym:} POMDP z QH dyskontowaniem
    \item \textbf{Multi-objective RL:} Równoważenie wielu kryteriów z różnymi preferencjami czasowymi
\end{itemize}

\vspace{0.5em}
\textbf{Interdyscyplinarne możliwości:}
\begin{itemize}
    \item Współpraca z ekonomistami behawioralnymi
    \item Badania kliniczne w psychologii uzależnień
    \item Aplikacje w teorii gier i projektowaniu mechanizmów
    \item Neuroekonomia -- połączenie z badaniami neuronaukowymi
\end{itemize}
\end{frame}

\begin{frame}[allowframebreaks]{Bibliografia}
\small
\begin{thebibliography}{99}

\bibitem{eshwar2025}
S. R. Eshwar.
\newblock Teaching Precommitted Agents: Model-Free Policy Evaluation and Control in Quasi-Hyperbolic Discounted MDPs.
\newblock In \emph{Proceedings of the 64th IEEE Conference on Decision and Control (CDC)}, 2025.

\bibitem{jaskiewicz2021}
A. Jaśkiewicz and A. S. Nowak.
\newblock Markov decision processes with quasi-hyperbolic discounting.
\newblock \emph{Finance and Stochastics}, 25(2):209--229, 2021.

\bibitem{eshwar2024}
S. R. Eshwar, M. Motwani, N. Roy, and G. Thoppe.
\newblock Reinforcement learning with quasi-hyperbolic discounting.
\newblock In \emph{IEEE Conference on Decision and Control (CDC)}, 2024.

\bibitem{borkar2008}
V. S. Borkar.
\newblock \emph{Stochastic Approximation: A Dynamical Systems Viewpoint}.
\newblock Cambridge University Press, 2008.

\bibitem{puterman2014}
M. L. Puterman.
\newblock \emph{Markov decision processes: discrete stochastic dynamic programming}.
\newblock John Wiley \& Sons, 2014.

\bibitem{Kallenberg2022}
L. C. M. Kallenberg.
\newblock \emph{Lecture Notes Markov Decision Processes}.
\newblock Leiden University, 2022.

\bibitem{sutton2018}
R. S. Sutton and A. G. Barto.
\newblock \emph{Reinforcement learning: An introduction}.
\newblock MIT press, 2nd edition, 2018.

\bibitem{laibson1997}
D. Laibson.
\newblock Golden eggs and hyperbolic discounting.
\newblock \emph{The Quarterly Journal of Economics}, 112(2):443--478, 1997.

\bibitem{bertsekas2019}
D. P. Bertsekas.
\newblock \emph{Reinforcement learning and optimal control}.
\newblock Athena Scientific, 2019.

\end{thebibliography}
\end{frame}

\begin{frame}[plain]
\begin{center}
{\Huge Dziękuję za uwagę}

\vspace{2em}
{\Large Pytania?}

\vspace{3em}
\normalsize
Michał Wrona\\
Politechnika Wrocławska\\
Wydział Matematyki\\[0.5em]
Promotor: prof. dr hab. inż. Anna Jaśkiewicz
\end{center}
\end{frame}
